% META: {"id": "TSPE-GEOESPACE-EX-044", "chapitre": "TSPE-GEOMETRIE-ESPACE", "type_objet": "exercice", "capacites": ["TSPE-GEOMETRIE-ESPACE-C14"], "capacites_codes": ["C14"], "methodes": ["M1"], "parcours": 1, "competences": ["calculer"], "duree_min": 8, "mode_creation": "ex_nihilo", "sources_inspiration": [], "fichier_tex": "chapitres/TSPE-GEOMETRIE-ESPACE/exercices/TSPE-GEOESPACE-EX-044.tex", "corrige_tex": "chapitres/TSPE-GEOMETRIE-ESPACE/corriges/TSPE-GEOESPACE-CO-044.tex", "status": "approved"}
% BEGIN-VERIFY
% A=(2,-1,0); u=(1,2,-1)
% t=2
% pt = tuple(A[i]+t*u[i] for i in range(3))
% assert pt == (4,3,-2)
% END-VERIFY
\begin{exercice}{TSPE-GEOESPACE-EX-002}{1}{10}
Soit $\mathcal{D}$ la droite passant par $A(2,-1,0)$ de vecteur directeur $\vec{u}\begin{pmatrix}1\\2\\-1\end{pmatrix}$.
\begin{enumerate}
  \item Ecrire une representation parametrique de $\mathcal{D}$.
  \item Le point $B(4,3,-2)$ appartient-il a $\mathcal{D}$ ?
\end{enumerate}
\end{exercice}
